% !TeX root = ../papers/nc-p3.tex
\section{三つの薄い色を持つ対称族の無印同型}
\label{sec:three-lines-symmetric-isomorphism}
\begin{theorem}\label{thm:three-lines-symmetric-isomorphism}
定理\ref{thm:klein-three-lines}の対称表示の二つの正則窓では、原則として四係数 $p_i$ の多重集合の一致が
添字保持同型の必要十分条件である。例外は二値がそれぞれ二回現れる場合だけである。
この場合の多重集合を $\{a,a,b,b\}$、$r=(a-b)/(a+b)$ と書く。
同じ型の二窓は $r'^2=r^2$ または $r'^2=r^2/(r^2-1)$ のとき、またそのときに限り同型である。
他の型との同型はない。
\end{theorem}
\begin{proof}[四つの二次式の束と自己随伴作用素]
系\ref{cor:symmetric-chart-isomorphism}より二次式の四次元空間 $L$ の合同を判定すればよい。
元の文字順で $p=(a,b,c,d)$ とすると、その行列は
\[
 \begin{pmatrix}
 ta&x&y&z\\ x&tb&-z&-y\\ y&-z&tc&-x\\ z&-y&-x&td
 \end{pmatrix}.
\]
$t=0$ の網を $N$ とする。一文字の符号を変え、四文字の正規直交Fourier基底に移すと、
$N$ は跡零の全対角行列、残る第四行列は対角成分が全て $1/4$ の四元群の巡回行列になる。
非対角成分は色ごとに一定で、4倍すると順序を除いて
$r=a+b-c-d$, $s=a+c-b-d$, $h=a+d-b-c$ となる。

全 $Q\in L$ について $T^tQ=QT$ を満たす作用素の空間を取る。$N$ の条件から $T$ は対角であり、
第四行列の非零な非対角成分の両端で対角値が一致する。従って $(r,s,h)$ の非零成分が
少なくとも二つ、一つ、零の場合に、それぞれ次元1、2、4である。これは合同不変量なので三者を分離する。
次元4なら $p_i$ は全て $1/4$ で一つの類である。
\end{proof}

\begin{proof}[自己随伴空間が一次元の場合の網の復元]
$\mathbb P N$ に階数2の行列が六つあり、跡零対角網の行列式は四直線に分かれる。
$t=1$ の階数2以下の点を調べる。$X=4x,Y=4y,Z=4z$ とする。相補的な二つの余因子の和を取ると
\[
 rYZ=(sh-r)X,\quad sXZ=(rh-s)Y,\quad hXY=(rs-h)Z .
\]
$ r,s,h$ の一つだけが零ならこれらは $X=Y=Z=0$ を強制し、元の対角行列は可逆である。
三つとも非零なら右辺の倍率を左の係数で割り $\alpha=sh/r-1,\beta=rh/s-1,\gamma=rs/h-1$ と置く。
$\alpha=0$ なら $r=sh$ のもとで $4p_i$ は $(1\pm s)(1\pm h)$ の四通りであり、
どれも非零だから $\beta\gamma\ne0$。従って $Y=Z=0$、二つの対角二行ブロックの行列式から
$X$ は高々二通りである。他の倍率が零でも同様である。

$\alpha\beta\gamma\ne0$ のとき、座標零の解は除かれるため
$X^2=\beta\gamma$, $Y^2=\alpha\gamma$, $Z^2=\alpha\beta$ である。
四符号の候補はアフィン独立である。他の平面が六節点（階数2の点）を含むなら、
そのうち三点がアフィン、三点が無限遠でなければならない。無限遠の三点を通る直線は
四直線の一つ $\pm X\pm Y\pm Z=0$ である。実際六点は四直線の二本ずつの交点で、他の三点直線はない。
同じ一次式が四符号点の三点で同じ値になると $X^2=Y^2=Z^2$、
従って $r^2=s^2=h^2$ となる。これは $p_i$ の三つが等しい場合である。
それ以外では六節点の平面は $\mathbb P N$ だけである。

この網の射影安定群はFourier基底の符号付き置換である。四本の行列式直線の各一般点の核が
四座標点を復元し、跡零を保つ対角変換の四平方は等しいからである。
第四行列の対角の跡を揃えた後、非零な相補二辺の成分が等しくなるには四符号の積が $+1$ でなければならない。
すると三つのFourier係数に対する作用は置換と偶数の符号反転であり、これは元の四つの $p_i$ の全置換に正確に対応する。
\end{proof}

\begin{proof}[五つの網が現れる場合]
三つが等しい配置を $(a,b,b,b)$, $a\ne b$ とする。Fourier基底を選び直すと第四行列は $bI+(a-b)J/4$ である。
束は座標四形式とその和の平方、合わせて五つの階数1二次式の張る五次元空間内で
\[
 (1,1,1,1,-16b/(a-b))
\]
を法線とする超平面である。五つの一次形式は射影枠をなし、どの四つも独立である。
法線のどの成分も零でない。五座標のうち二つだけが非零となる十点が階数2であり、これで全てである。
三つか四つだけが非零ならその個数が階数であり、五つ全部が非零なら可逆な対角行列に階数1を加えても階数3以上となる。

法線の各成分で五座標をそれぞれ倍率変更すれば、これらの十点は成分和零の空間の $e_i-e_j$ である。
平面上に入る対は平面の方程式の五係数が等しい対に対応するため、六点を持つのは四係数が等しい五平面だけである。
各平面の網は対応する四つの独立な一次形式の平方で対角化される。
四本の行列式直線の核から双対の四座標形式を復元すると、その平方全体の張る空間も復元される。
これら五つの空間の和が元の五次元空間であり、その階数1点はちょうど五つである。
ゆえにこの射影枠は無印で復元される。五一次形式を和が零になるように規格化すれば、その置換に伴う倍率は
共通スカラーだけであり、法線の係数は置換と共通倍率まで復元される。
四つが1の多重集合は最後の値を一意に決める（五つとも1の場合を含む）。$a+3b=1$ より元の係数も決まる。
\end{proof}

\begin{proof}[二つの二重値]
自己随伴空間が二次元なら、四係数は二値の二対であり、二つの原始冪等元が二平面の無順序対を復元する。
第四行列の対角を1に規格化し、非零な二本の相補辺の成分を $r$ とする。
$r\ne0,\pm1$ である。一般の行列の対角を $d_0,\ldots,d_3$ とすると $t=(\sum d_i)/4$、
二本の非対角成分は $rt$ である。片方の二平面だけに支えられる二次式は各々 $x^2-y^2$ の直線であり、
その等方二直線が決まる。これらの式を割った後の各二平面での座標を、等方二直線での値
$U_\pm=d_0+d_1\pm2rt$, $V_\pm=d_2+d_3\pm2rt$ に取る。
二平面間のグラフ行列は
\[
 \binom{V_+}{V_-}=\frac1r\begin{pmatrix}1&-1-r\\1-r&-1\end{pmatrix}\binom{U_+}{U_-}.
\]
各等方対を保つ変換は二行・二列を別々に置換・倍率変更できる。値に現れる倍率は平方だが代数閉体上で任意に取れる。
逆にこれで商のグラフを合わせれば、その逆像である元の束も一致する。二平面の交換は逆行列を取る。
全成分非零の二行二列行列の行列倍率商は両対角積の比で決まり、
ここでは $1/(1-r^2)$ とその逆だけが同一視される。これが定理の二条件である。
\path{computations/verify_three_lines_symmetric_isomorphism.py} がFourier変換、余因子、五平面と二平面のグラフを検算する。
\end{proof}
